\chapter{AgentOS Kernel 的系统结构}
\chaptervoice{这一章，我们围绕“AgentOS Kernel 的系统结构”做一件克制的事：把直觉压缩为状态、关系、时间尺度和可检验后果。它在全书中的任务，是说明认知内核通过身体进入现实。}

\section*{理论来路与依赖接口}
本章的直接思想来源是原文第 353614—368709 行的 KRA 与双轴多尺度 AgentOS 收束稿。我们采用原文较晚形成的定义来整理较早讨论，不把对话中的试探性比喻与最终模型并列成互相竞争的“定律”。已有理论锚点主要是具身认知、分层控制、MPC、状态估计与安全包络（模型预测控制、Kalman 状态估计与具身认知）。

本章使用上一章“Linux 之上的 AgentOS”已经得到的对象与边界；本章产出将供下一章“Agent Kernel—KRA—Body Runtime 完整数据流”使用。若后文术语偶尔出现，只作为路线提示，不参与本章推导。

\section*{从哪里出发}
我们已经拥有上一章给出的概念，现在只向前走一步。我会先把对象放进闭环，再讨论它的数学形式，最后检查工程边界。读者不必预先相信本章术语；只需暂时接受一个方法论约定：智能结构的意义来自它在现实反馈中的作用，而不是来自名称本身。

令系统在观察窗口 $[t,t+T]$ 内的状态轨迹为 $x_{[t,t+T]}$，环境扰动为 $d$，行动为 $u$。本章讨论的对象若确实有效，就应改变条件分布 $p(x_{t+T}\mid x_t,u,d)$，或改变达到同一结果所需的代价。这个起点把讨论锚定在具身认知、分层控制、MPC、状态估计与安全包络上。
\section{Persistent Subject Manager}
\begin{narration}
现在把镜头移到“Persistent Subject Manager”。我们只沿本章已经取得的概念向前一步：先确定它在闭环里的角色，再检查它的尺度与失败边界。
\end{narration}

本节把“Persistent Subject Manager”落实为在本章闭环中具有明确输入、输出、状态和时间尺度的功能关系。这个定义不是说“任何相似现象都算”，而是要求我们同时给出对象域、允许的干预以及观察窗口。对于主题“AgentOS Kernel 的系统结构”而言，它承担的是认知内核通过身体进入现实中的一个可辨识环节。

把这一关系写成最小数学骨架，可得

\begin{equation}z_{t+1}=F(z_t,i_t;\alpha),\qquad o_t=H(z_t).\end{equation}

式中的符号只承担结构角色，具体参数仍需由任务和身体数据辨识。我们首先检查可观测量，再检查控制量，最后检查比它更慢的约束。这样的顺序来自具身认知、分层控制、MPC、状态估计与安全包络，但本书对Persistent Subject Manager的命名和组合仍是需要实验验证的建模提案。

这一小节的反例同样重要：若不能通过干预改变轨迹分布或资源代价，该关系在当前尺度上不可辨识。因此评价它不能只看一次任务结果，而要比较干预前后的轨迹、恢复时间、维护成本和风险。如果改变该机制而所有允许观察下的分布都不变，那么它至少在当前实验族中是冗余描述。

\begin{thought}
对“Persistent Subject Manager”做一次消融：冻结它、删除它，或只允许它以相邻时间尺度交换残差。哪些现象立即消失，哪些只在长期出现？若答案无法区分三种干预，我们还没有找到正确的状态变量。
\end{thought}

最后，把结论交还现实：本节只建立功能层关系，不由此推出特定脑区实现、主观意识或宇宙目的。可测状态、干预后果和明确失败条件，才是从原文直觉走向可检验理论的桥。

\section{h Cognitive Bus}
\begin{narration}
现在把镜头移到“h Cognitive Bus”。我们只沿本章已经取得的概念向前一步：先确定它在闭环里的角色，再检查它的尺度与失败边界。
\end{narration}

本节把“h Cognitive Bus”落实为在本章闭环中具有明确输入、输出、状态和时间尺度的功能关系。这个定义不是说“任何相似现象都算”，而是要求我们同时给出对象域、允许的干预以及观察窗口。对于主题“AgentOS Kernel 的系统结构”而言，它承担的是认知内核通过身体进入现实中的一个可辨识环节。

把这一关系写成最小数学骨架，可得

\begin{equation}z_{t+1}=F(z_t,i_t;\alpha),\qquad o_t=H(z_t).\end{equation}

式中的符号只承担结构角色，具体参数仍需由任务和身体数据辨识。我们首先检查可观测量，再检查控制量，最后检查比它更慢的约束。这样的顺序来自具身认知、分层控制、MPC、状态估计与安全包络，但本书对h Cognitive Bus的命名和组合仍是需要实验验证的建模提案。

这一小节的反例同样重要：若不能通过干预改变轨迹分布或资源代价，该关系在当前尺度上不可辨识。因此评价它不能只看一次任务结果，而要比较干预前后的轨迹、恢复时间、维护成本和风险。如果改变该机制而所有允许观察下的分布都不变，那么它至少在当前实验族中是冗余描述。

\begin{thought}
对“h Cognitive Bus”做一次消融：冻结它、删除它，或只允许它以相邻时间尺度交换残差。哪些现象立即消失，哪些只在长期出现？若答案无法区分三种干预，我们还没有找到正确的状态变量。
\end{thought}

最后，把结论交还现实：本节只建立功能层关系，不由此推出特定脑区实现、主观意识或宇宙目的。可测状态、干预后果和明确失败条件，才是从原文直觉走向可检验理论的桥。

\section{Lifecycle Memory Manager}
\begin{narration}
现在把镜头移到“Lifecycle Memory Manager”。我们只沿本章已经取得的概念向前一步：先确定它在闭环里的角色，再检查它的尺度与失败边界。
\end{narration}

本节把“Lifecycle Memory Manager”落实为在本章闭环中具有明确输入、输出、状态和时间尺度的功能关系。这个定义不是说“任何相似现象都算”，而是要求我们同时给出对象域、允许的干预以及观察窗口。对于主题“AgentOS Kernel 的系统结构”而言，它承担的是认知内核通过身体进入现实中的一个可辨识环节。

把这一关系写成最小数学骨架，可得

\begin{equation}z_{t+1}=F(z_t,i_t;\alpha),\qquad o_t=H(z_t).\end{equation}

式中的符号只承担结构角色，具体参数仍需由任务和身体数据辨识。我们首先检查可观测量，再检查控制量，最后检查比它更慢的约束。这样的顺序来自具身认知、分层控制、MPC、状态估计与安全包络，但本书对Lifecycle Memory Manager的命名和组合仍是需要实验验证的建模提案。

这一小节的反例同样重要：若不能通过干预改变轨迹分布或资源代价，该关系在当前尺度上不可辨识。因此评价它不能只看一次任务结果，而要比较干预前后的轨迹、恢复时间、维护成本和风险。如果改变该机制而所有允许观察下的分布都不变，那么它至少在当前实验族中是冗余描述。

\begin{thought}
对“Lifecycle Memory Manager”做一次消融：冻结它、删除它，或只允许它以相邻时间尺度交换残差。哪些现象立即消失，哪些只在长期出现？若答案无法区分三种干预，我们还没有找到正确的状态变量。
\end{thought}

最后，把结论交还现实：本节只建立功能层关系，不由此推出特定脑区实现、主观意识或宇宙目的。可测状态、干预后果和明确失败条件，才是从原文直觉走向可检验理论的桥。

\section{Scheduler}
\begin{narration}
现在把镜头移到“Scheduler”。我们只沿本章已经取得的概念向前一步：先确定它在闭环里的角色，再检查它的尺度与失败边界。
\end{narration}

本节把“Scheduler”落实为在有限显式工作表面上执行复杂度整形、资源竞争、动态准入与可恢复外移的治理机制。这个定义不是说“任何相似现象都算”，而是要求我们同时给出对象域、允许的干预以及观察窗口。对于主题“AgentOS Kernel 的系统结构”而言，它承担的是认知内核通过身体进入现实中的一个可辨识环节。

把这一关系写成最小数学骨架，可得

\begin{equation}\sum_i c_iq_i\leq W,\quad q_i\in\{0,1\}.\end{equation}

式中的符号只承担结构角色，具体参数仍需由任务和身体数据辨识。我们首先检查可观测量，再检查控制量，最后检查比它更慢的约束。这样的顺序来自具身认知、分层控制、MPC、状态估计与安全包络，但本书对Scheduler的命名和组合仍是需要实验验证的建模提案。

这一小节的反例同样重要：把压缩、调度、边界和卸载并成同一模块会使权限与失败责任不清。因此评价它不能只看一次任务结果，而要比较干预前后的轨迹、恢复时间、维护成本和风险。如果改变该机制而所有允许观察下的分布都不变，那么它至少在当前实验族中是冗余描述。

\begin{thought}
对“Scheduler”做一次消融：冻结它、删除它，或只允许它以相邻时间尺度交换残差。哪些现象立即消失，哪些只在长期出现？若答案无法区分三种干预，我们还没有找到正确的状态变量。
\end{thought}

最后，把结论交还现实：本节只建立功能层关系，不由此推出特定脑区实现、主观意识或宇宙目的。可测状态、干预后果和明确失败条件，才是从原文直觉走向可检验理论的桥。

\section{SPC}
\begin{narration}
现在把镜头移到“SPC”。我们只沿本章已经取得的概念向前一步：先确定它在闭环里的角色，再检查它的尺度与失败边界。
\end{narration}

本节把“SPC”落实为从分布式内部状态压缩出带置信度的后验自状态估计器。这个定义不是说“任何相似现象都算”，而是要求我们同时给出对象域、允许的干预以及观察窗口。对于主题“AgentOS Kernel 的系统结构”而言，它承担的是认知内核通过身体进入现实中的一个可辨识环节。

把这一关系写成最小数学骨架，可得

\begin{equation}\hat z_t,\Sigma_t=\mathsf{SPC}(h_{[t-\Delta,t]},a_t).\end{equation}

式中的符号只承担结构角色，具体参数仍需由任务和身体数据辨识。我们首先检查可观测量，再检查控制量，最后检查比它更慢的约束。这样的顺序来自具身认知、分层控制、MPC、状态估计与安全包络，但本书对SPC的命名和组合仍是需要实验验证的建模提案。

这一小节的反例同样重要：把观察器输出当作透明自知会隐藏偏差、方差与延迟。因此评价它不能只看一次任务结果，而要比较干预前后的轨迹、恢复时间、维护成本和风险。如果改变该机制而所有允许观察下的分布都不变，那么它至少在当前实验族中是冗余描述。

\begin{thought}
对“SPC”做一次消融：冻结它、删除它，或只允许它以相邻时间尺度交换残差。哪些现象立即消失，哪些只在长期出现？若答案无法区分三种干预，我们还没有找到正确的状态变量。
\end{thought}

最后，把结论交还现实：本节只建立功能层关系，不由此推出特定脑区实现、主观意识或宇宙目的。可测状态、干预后果和明确失败条件，才是从原文直觉走向可检验理论的桥。

\section{AHC}
\begin{narration}
现在把镜头移到“AHC”。我们只沿本章已经取得的概念向前一步：先确定它在闭环里的角色，再检查它的尺度与失败边界。
\end{narration}

本节把“AHC”落实为对威胁、唤醒、能量、疲劳、信任等信号进行积分、衰减和饱和的连续调制场。这个定义不是说“任何相似现象都算”，而是要求我们同时给出对象域、允许的干预以及观察窗口。对于主题“AgentOS Kernel 的系统结构”而言，它承担的是认知内核通过身体进入现实中的一个可辨识环节。

把这一关系写成最小数学骨架，可得

\begin{equation}\dot a=-\Lambda(a-a_0)+Bs,\quad a\in[a_-,a_+].\end{equation}

式中的符号只承担结构角色，具体参数仍需由任务和身体数据辨识。我们首先检查可观测量，再检查控制量，最后检查比它更慢的约束。这样的顺序来自具身认知、分层控制、MPC、状态估计与安全包络，但本书对AHC的命名和组合仍是需要实验验证的建模提案。

这一小节的反例同样重要：拟人化命名若没有功能通道会把工程状态误写成主观感受。因此评价它不能只看一次任务结果，而要比较干预前后的轨迹、恢复时间、维护成本和风险。如果改变该机制而所有允许观察下的分布都不变，那么它至少在当前实验族中是冗余描述。

\begin{thought}
对“AHC”做一次消融：冻结它、删除它，或只允许它以相邻时间尺度交换残差。哪些现象立即消失，哪些只在长期出现？若答案无法区分三种干预，我们还没有找到正确的状态变量。
\end{thought}

最后，把结论交还现实：本节只建立功能层关系，不由此推出特定脑区实现、主观意识或宇宙目的。可测状态、干预后果和明确失败条件，才是从原文直觉走向可检验理论的桥。

\section{TCE}
\begin{narration}
现在把镜头移到“TCE”。我们只沿本章已经取得的概念向前一步：先确定它在闭环里的角色，再检查它的尺度与失败边界。
\end{narration}

本节把“TCE”落实为依据自状态估计调节探索、注意范围、搜索深度、步长与保守度的受限控制器。这个定义不是说“任何相似现象都算”，而是要求我们同时给出对象域、允许的干预以及观察窗口。对于主题“AgentOS Kernel 的系统结构”而言，它承担的是认知内核通过身体进入现实中的一个可辨识环节。

把这一关系写成最小数学骨架，可得

\begin{equation}u_c=\operatorname{sat}_{\mathcal U}(K(r-\hat z)).\end{equation}

式中的符号只承担结构角色，具体参数仍需由任务和身体数据辨识。我们首先检查可观测量，再检查控制量，最后检查比它更慢的约束。这样的顺序来自具身认知、分层控制、MPC、状态估计与安全包络，但本书对TCE的命名和组合仍是需要实验验证的建模提案。

这一小节的反例同样重要：增益过高和观察延迟叠加会造成过冲与认知振荡。因此评价它不能只看一次任务结果，而要比较干预前后的轨迹、恢复时间、维护成本和风险。如果改变该机制而所有允许观察下的分布都不变，那么它至少在当前实验族中是冗余描述。

\begin{thought}
对“TCE”做一次消融：冻结它、删除它，或只允许它以相邻时间尺度交换残差。哪些现象立即消失，哪些只在长期出现？若答案无法区分三种干预，我们还没有找到正确的状态变量。
\end{thought}

最后，把结论交还现实：本节只建立功能层关系，不由此推出特定脑区实现、主观意识或宇宙目的。可测状态、干预后果和明确失败条件，才是从原文直觉走向可检验理论的桥。

\section{CCM}
\begin{narration}
现在把镜头移到“CCM”。我们只沿本章已经取得的概念向前一步：先确定它在闭环里的角色，再检查它的尺度与失败边界。
\end{narration}

本节把“CCM”落实为把显式工作负担分解、去重、折叠或重构到容量约束内的复杂度整形器。这个定义不是说“任何相似现象都算”，而是要求我们同时给出对象域、允许的干预以及观察窗口。对于主题“AgentOS Kernel 的系统结构”而言，它承担的是认知内核通过身体进入现实中的一个可辨识环节。

把这一关系写成最小数学骨架，可得

\begin{equation}C_{\rm run}(t)\leq C_{\max}(W).\end{equation}

式中的符号只承担结构角色，具体参数仍需由任务和身体数据辨识。我们首先检查可观测量，再检查控制量，最后检查比它更慢的约束。这样的顺序来自具身认知、分层控制、MPC、状态估计与安全包络，但本书对CCM的命名和组合仍是需要实验验证的建模提案。

这一小节的反例同样重要：若 CCM 同时决定权限和任务价值，复杂度治理会越权替代调度。因此评价它不能只看一次任务结果，而要比较干预前后的轨迹、恢复时间、维护成本和风险。如果改变该机制而所有允许观察下的分布都不变，那么它至少在当前实验族中是冗余描述。

\begin{thought}
对“CCM”做一次消融：冻结它、删除它，或只允许它以相邻时间尺度交换残差。哪些现象立即消失，哪些只在长期出现？若答案无法区分三种干预，我们还没有找到正确的状态变量。
\end{thought}

最后，把结论交还现实：本节只建立功能层关系，不由此推出特定脑区实现、主观意识或宇宙目的。可测状态、干预后果和明确失败条件，才是从原文直觉走向可检验理论的桥。

\section{Boundary}
\begin{narration}
现在把镜头移到“Boundary”。我们只沿本章已经取得的概念向前一步：先确定它在闭环里的角色，再检查它的尺度与失败边界。
\end{narration}

本节把“Boundary”落实为在有限显式工作表面上执行复杂度整形、资源竞争、动态准入与可恢复外移的治理机制。这个定义不是说“任何相似现象都算”，而是要求我们同时给出对象域、允许的干预以及观察窗口。对于主题“AgentOS Kernel 的系统结构”而言，它承担的是认知内核通过身体进入现实中的一个可辨识环节。

把这一关系写成最小数学骨架，可得

\begin{equation}\sum_i c_iq_i\leq W,\quad q_i\in\{0,1\}.\end{equation}

式中的符号只承担结构角色，具体参数仍需由任务和身体数据辨识。我们首先检查可观测量，再检查控制量，最后检查比它更慢的约束。这样的顺序来自具身认知、分层控制、MPC、状态估计与安全包络，但本书对Boundary的命名和组合仍是需要实验验证的建模提案。

这一小节的反例同样重要：把压缩、调度、边界和卸载并成同一模块会使权限与失败责任不清。因此评价它不能只看一次任务结果，而要比较干预前后的轨迹、恢复时间、维护成本和风险。如果改变该机制而所有允许观察下的分布都不变，那么它至少在当前实验族中是冗余描述。

\begin{thought}
对“Boundary”做一次消融：冻结它、删除它，或只允许它以相邻时间尺度交换残差。哪些现象立即消失，哪些只在长期出现？若答案无法区分三种干预，我们还没有找到正确的状态变量。
\end{thought}

最后，把结论交还现实：本节只建立功能层关系，不由此推出特定脑区实现、主观意识或宇宙目的。可测状态、干预后果和明确失败条件，才是从原文直觉走向可检验理论的桥。

\section{Offloading Manager}
\begin{narration}
现在把镜头移到“Offloading Manager”。我们只沿本章已经取得的概念向前一步：先确定它在闭环里的角色，再检查它的尺度与失败边界。
\end{narration}

本节把“Offloading Manager”落实为在有限显式工作表面上执行复杂度整形、资源竞争、动态准入与可恢复外移的治理机制。这个定义不是说“任何相似现象都算”，而是要求我们同时给出对象域、允许的干预以及观察窗口。对于主题“AgentOS Kernel 的系统结构”而言，它承担的是认知内核通过身体进入现实中的一个可辨识环节。

把这一关系写成最小数学骨架，可得

\begin{equation}\sum_i c_iq_i\leq W,\quad q_i\in\{0,1\}.\end{equation}

式中的符号只承担结构角色，具体参数仍需由任务和身体数据辨识。我们首先检查可观测量，再检查控制量，最后检查比它更慢的约束。这样的顺序来自具身认知、分层控制、MPC、状态估计与安全包络，但本书对Offloading Manager的命名和组合仍是需要实验验证的建模提案。

这一小节的反例同样重要：把压缩、调度、边界和卸载并成同一模块会使权限与失败责任不清。因此评价它不能只看一次任务结果，而要比较干预前后的轨迹、恢复时间、维护成本和风险。如果改变该机制而所有允许观察下的分布都不变，那么它至少在当前实验族中是冗余描述。

\begin{thought}
对“Offloading Manager”做一次消融：冻结它、删除它，或只允许它以相邻时间尺度交换残差。哪些现象立即消失，哪些只在长期出现？若答案无法区分三种干预，我们还没有找到正确的状态变量。
\end{thought}

最后，把结论交还现实：本节只建立功能层关系，不由此推出特定脑区实现、主观意识或宇宙目的。可测状态、干预后果和明确失败条件，才是从原文直觉走向可检验理论的桥。

\section{Goal/Task Runtime}
\begin{narration}
现在把镜头移到“Goal/Task Runtime”。我们只沿本章已经取得的概念向前一步：先确定它在闭环里的角色，再检查它的尺度与失败边界。
\end{narration}

本节把“Goal/Task Runtime”落实为在本章闭环中具有明确输入、输出、状态和时间尺度的功能关系。这个定义不是说“任何相似现象都算”，而是要求我们同时给出对象域、允许的干预以及观察窗口。对于主题“AgentOS Kernel 的系统结构”而言，它承担的是认知内核通过身体进入现实中的一个可辨识环节。

把这一关系写成最小数学骨架，可得

\begin{equation}z_{t+1}=F(z_t,i_t;\alpha),\qquad o_t=H(z_t).\end{equation}

式中的符号只承担结构角色，具体参数仍需由任务和身体数据辨识。我们首先检查可观测量，再检查控制量，最后检查比它更慢的约束。这样的顺序来自具身认知、分层控制、MPC、状态估计与安全包络，但本书对Goal/Task Runtime的命名和组合仍是需要实验验证的建模提案。

这一小节的反例同样重要：若不能通过干预改变轨迹分布或资源代价，该关系在当前尺度上不可辨识。因此评价它不能只看一次任务结果，而要比较干预前后的轨迹、恢复时间、维护成本和风险。如果改变该机制而所有允许观察下的分布都不变，那么它至少在当前实验族中是冗余描述。

\begin{thought}
对“Goal/Task Runtime”做一次消融：冻结它、删除它，或只允许它以相邻时间尺度交换残差。哪些现象立即消失，哪些只在长期出现？若答案无法区分三种干预，我们还没有找到正确的状态变量。
\end{thought}

最后，把结论交还现实：本节只建立功能层关系，不由此推出特定脑区实现、主观意识或宇宙目的。可测状态、干预后果和明确失败条件，才是从原文直觉走向可检验理论的桥。

\section{Skill Runtime}
\begin{narration}
现在把镜头移到“Skill Runtime”。我们只沿本章已经取得的概念向前一步：先确定它在闭环里的角色，再检查它的尺度与失败边界。
\end{narration}

本节把“Skill Runtime”落实为在本章闭环中具有明确输入、输出、状态和时间尺度的功能关系。这个定义不是说“任何相似现象都算”，而是要求我们同时给出对象域、允许的干预以及观察窗口。对于主题“AgentOS Kernel 的系统结构”而言，它承担的是认知内核通过身体进入现实中的一个可辨识环节。

把这一关系写成最小数学骨架，可得

\begin{equation}z_{t+1}=F(z_t,i_t;\alpha),\qquad o_t=H(z_t).\end{equation}

式中的符号只承担结构角色，具体参数仍需由任务和身体数据辨识。我们首先检查可观测量，再检查控制量，最后检查比它更慢的约束。这样的顺序来自具身认知、分层控制、MPC、状态估计与安全包络，但本书对Skill Runtime的命名和组合仍是需要实验验证的建模提案。

这一小节的反例同样重要：若不能通过干预改变轨迹分布或资源代价，该关系在当前尺度上不可辨识。因此评价它不能只看一次任务结果，而要比较干预前后的轨迹、恢复时间、维护成本和风险。如果改变该机制而所有允许观察下的分布都不变，那么它至少在当前实验族中是冗余描述。

\begin{thought}
对“Skill Runtime”做一次消融：冻结它、删除它，或只允许它以相邻时间尺度交换残差。哪些现象立即消失，哪些只在长期出现？若答案无法区分三种干预，我们还没有找到正确的状态变量。
\end{thought}

最后，把结论交还现实：本节只建立功能层关系，不由此推出特定脑区实现、主观意识或宇宙目的。可测状态、干预后果和明确失败条件，才是从原文直觉走向可检验理论的桥。

\section{External Memory Manager}
\begin{narration}
现在把镜头移到“External Memory Manager”。我们只沿本章已经取得的概念向前一步：先确定它在闭环里的角色，再检查它的尺度与失败边界。
\end{narration}

本节把“External Memory Manager”落实为在本章闭环中具有明确输入、输出、状态和时间尺度的功能关系。这个定义不是说“任何相似现象都算”，而是要求我们同时给出对象域、允许的干预以及观察窗口。对于主题“AgentOS Kernel 的系统结构”而言，它承担的是认知内核通过身体进入现实中的一个可辨识环节。

把这一关系写成最小数学骨架，可得

\begin{equation}z_{t+1}=F(z_t,i_t;\alpha),\qquad o_t=H(z_t).\end{equation}

式中的符号只承担结构角色，具体参数仍需由任务和身体数据辨识。我们首先检查可观测量，再检查控制量，最后检查比它更慢的约束。这样的顺序来自具身认知、分层控制、MPC、状态估计与安全包络，但本书对External Memory Manager的命名和组合仍是需要实验验证的建模提案。

这一小节的反例同样重要：若不能通过干预改变轨迹分布或资源代价，该关系在当前尺度上不可辨识。因此评价它不能只看一次任务结果，而要比较干预前后的轨迹、恢复时间、维护成本和风险。如果改变该机制而所有允许观察下的分布都不变，那么它至少在当前实验族中是冗余描述。

\begin{thought}
对“External Memory Manager”做一次消融：冻结它、删除它，或只允许它以相邻时间尺度交换残差。哪些现象立即消失，哪些只在长期出现？若答案无法区分三种干预，我们还没有找到正确的状态变量。
\end{thought}

最后，把结论交还现实：本节只建立功能层关系，不由此推出特定脑区实现、主观意识或宇宙目的。可测状态、干预后果和明确失败条件，才是从原文直觉走向可检验理论的桥。

\section{Checkpoint / Recovery}
\begin{narration}
现在把镜头移到“Checkpoint / Recovery”。我们只沿本章已经取得的概念向前一步：先确定它在闭环里的角色，再检查它的尺度与失败边界。
\end{narration}

本节把“Checkpoint / Recovery”落实为连续认知演化与离散目标、权限、技能和恢复事件组成的混杂系统。这个定义不是说“任何相似现象都算”，而是要求我们同时给出对象域、允许的干预以及观察窗口。对于主题“AgentOS Kernel 的系统结构”而言，它承担的是认知内核通过身体进入现实中的一个可辨识环节。

把这一关系写成最小数学骨架，可得

\begin{equation}\dot x=f_q(x),\quad x^+=R_e(x),\ e\in\mathcal E.\end{equation}

式中的符号只承担结构角色，具体参数仍需由任务和身体数据辨识。我们首先检查可观测量，再检查控制量，最后检查比它更慢的约束。这样的顺序来自具身认知、分层控制、MPC、状态估计与安全包络，但本书对Checkpoint / Recovery的命名和组合仍是需要实验验证的建模提案。

这一小节的反例同样重要：纯状态机丢失连续积累，纯连续系统则无法表达治理权限。因此评价它不能只看一次任务结果，而要比较干预前后的轨迹、恢复时间、维护成本和风险。如果改变该机制而所有允许观察下的分布都不变，那么它至少在当前实验族中是冗余描述。

\begin{thought}
对“Checkpoint / Recovery”做一次消融：冻结它、删除它，或只允许它以相邻时间尺度交换残差。哪些现象立即消失，哪些只在长期出现？若答案无法区分三种干预，我们还没有找到正确的状态变量。
\end{thought}

最后，把结论交还现实：本节只建立功能层关系，不由此推出特定脑区实现、主观意识或宇宙目的。可测状态、干预后果和明确失败条件，才是从原文直觉走向可检验理论的桥。

\section{Governance Interface}
\begin{narration}
现在把镜头移到“Governance Interface”。我们只沿本章已经取得的概念向前一步：先确定它在闭环里的角色，再检查它的尺度与失败边界。
\end{narration}

本节把“Governance Interface”落实为在本章闭环中具有明确输入、输出、状态和时间尺度的功能关系。这个定义不是说“任何相似现象都算”，而是要求我们同时给出对象域、允许的干预以及观察窗口。对于主题“AgentOS Kernel 的系统结构”而言，它承担的是认知内核通过身体进入现实中的一个可辨识环节。

把这一关系写成最小数学骨架，可得

\begin{equation}z_{t+1}=F(z_t,i_t;\alpha),\qquad o_t=H(z_t).\end{equation}

式中的符号只承担结构角色，具体参数仍需由任务和身体数据辨识。我们首先检查可观测量，再检查控制量，最后检查比它更慢的约束。这样的顺序来自具身认知、分层控制、MPC、状态估计与安全包络，但本书对Governance Interface的命名和组合仍是需要实验验证的建模提案。

这一小节的反例同样重要：若不能通过干预改变轨迹分布或资源代价，该关系在当前尺度上不可辨识。因此评价它不能只看一次任务结果，而要比较干预前后的轨迹、恢复时间、维护成本和风险。如果改变该机制而所有允许观察下的分布都不变，那么它至少在当前实验族中是冗余描述。

\begin{thought}
对“Governance Interface”做一次消融：冻结它、删除它，或只允许它以相邻时间尺度交换残差。哪些现象立即消失，哪些只在长期出现？若答案无法区分三种干预，我们还没有找到正确的状态变量。
\end{thought}

最后，把结论交还现实：本节只建立功能层关系，不由此推出特定脑区实现、主观意识或宇宙目的。可测状态、干预后果和明确失败条件，才是从原文直觉走向可检验理论的桥。

\section*{本章收束：把名词还原为闭环}
现在回头看“AgentOS Kernel 的系统结构”，最重要的不是记住缩写，而是说清本章对象怎样改变系统的状态、代价或可恢复性。下一章“Agent Kernel—KRA—Body Runtime 完整数据流”只使用这里已经给出的结果，不反过来为本章提供前提。

\begin{bookproposition}
若一个候选机制既不改变可观测轨迹分布，也不改变资源、风险或恢复代价，那么在当前观察尺度上，它不是一个可辨识的功能结构。
\end{bookproposition}
\begin{proofidea}
功能等价由可干预后果定义。若所有允许干预下的输出分布与代价均不变，则该机制相对于当前实验族不可辨识；要么扩大观察尺度，要么删除这一冗余描述。
\end{proofidea}

\begin{readercheck}
请尝试回答：本章对象的状态是什么？它的最快和最慢时间常数是什么？什么观测能证伪它？它失败时，残差应向哪一层升级？
\end{readercheck}
\cleardoublepage
